\documentclass[12pt,a4paper]{article}

% ============================================================
% PetroLúmen - Memorando Estratégico e Técnico (Versão 1.0)
% Compilação recomendada: pdfLaTeX (executar duas vezes para atualizar o sumário)
% ============================================================

% ------------------------------------------------------------
% PACOTES
% ------------------------------------------------------------
\usepackage[utf8]{inputenc}
\usepackage[T1]{fontenc}
\usepackage[portuguese]{babel}
\usepackage{lmodern}
\usepackage{microtype}
\usepackage[a4paper,top=2.7cm,bottom=2.5cm,left=2.6cm,right=2.4cm,headheight=15pt]{geometry}
\usepackage{booktabs}
\usepackage{array}
\usepackage{tabularx}
\usepackage{longtable}
\usepackage{makecell}
\usepackage{colortbl}
\usepackage{xcolor}
\usepackage{enumitem}
\usepackage{eurosym}
\usepackage{parskip}
\usepackage{titlesec}
\usepackage{fancyhdr}
\usepackage{graphicx}
\usepackage{amsmath,amssymb}
\usepackage{caption}
\usepackage{float}
\usepackage{hyperref}

% ------------------------------------------------------------
% PALETA DE CORES INSTITUCIONAL (Baseado no audit report)
% ------------------------------------------------------------
\definecolor{primary}{RGB}{12,45,85}      % Azul escuro
\definecolor{secondary}{RGB}{0,120,110}   % Verde tecnológico
\definecolor{accent}{RGB}{200,80,20}      % Laranja destaque
\definecolor{lightgray}{RGB}{245,247,249}
\definecolor{tableheader}{RGB}{12,45,85}

% ------------------------------------------------------------
% CONFIGURAÇÕES DE LINKS E TÍTULOS
% ------------------------------------------------------------
\hypersetup{
  colorlinks=true, linkcolor=primary, urlcolor=secondary, citecolor=primary,
  pdftitle={PetroLúmen - Plataforma de Engenharia de Reservatórios},
  pdfauthor={PetroLúmen Team},
  pdfsubject={Auditoria de Arquitetura e Validação Técnica de Software de Reservatórios},
  pdfcreator={LaTeX}
}

\titleformat{\section}{\Large\bfseries\color{primary}}{\thesection}{0.8em}{}
\titleformat{\subsection}{\large\bfseries\color{secondary}}{\thesubsection}{0.8em}{}
\titleformat{\subsubsection}{\normalsize\bfseries\color{accent}}{\thesubsubsection}{0.8em}{}
\setcounter{secnumdepth}{3}
\setcounter{tocdepth}{2}
\setlength{\emergencystretch}{3em}
\setlength{\tabcolsep}{4.5pt}
\renewcommand{\arraystretch}{1.15}
\setlist[itemize]{leftmargin=*,topsep=3pt,itemsep=2pt}
\setlist[enumerate]{leftmargin=*,topsep=3pt,itemsep=3pt}

\pagestyle{fancy}
\fancyhf{}
\fancyhead[L]{\small\textcolor{primary}{PetroLúmen}}
\fancyhead[R]{\small\textcolor{primary}{Memorando Estratégico e Técnico | Rev. 1.0}}
\fancyfoot[C]{\thepage}
\renewcommand{\headrulewidth}{0.4pt}
\renewcommand{\footrulewidth}{0pt}
\fancypagestyle{plain}{
  \fancyhf{}\fancyfoot[C]{\thepage}\renewcommand{\headrulewidth}{0pt}
}
\captionsetup{font=small,labelfont=bf,labelsep=period}

% ------------------------------------------------------------
% COMANDOS PERSONALIZADOS
% ------------------------------------------------------------
\newcommand{\usd}{US\$}
\newcommand{\statusvalidar}{\textcolor{accent}{\textbf{A validar}}}

% ============================================================
% DOCUMENTO
% ============================================================
\sloppy
\tolerance=3000
\hbadness=10000
\vbadness=10000
\raggedbottom

\begin{document}

% ------------------------------------------------------------
% CAPA
% ------------------------------------------------------------
\begin{titlepage}
  \thispagestyle{empty}
  \centering
  \vspace*{1.8cm}
  {\Huge\bfseries\color{primary} PetroLúmen\par}
  \vspace{0.6cm}
  {\Large\color{secondary} Plataforma Avançada de Engenharia de Reservatórios\par}
  \vspace{1.4cm}
  \rule{\textwidth}{1.4pt}\\[0.45cm]
  {\large\bfseries Memorando Estratégico, Arquitetura Técnica e Validação de Escalabilidade\par}
  \rule{\textwidth}{1.4pt}\\[1.5cm]
  \begin{minipage}{0.86\textwidth}
    \centering\large
    \textit{``Análise computacional de alta performance, com isolamento de sessão e motores de simulação determinísticos.''}
  \end{minipage}
  \vfill
  \begin{tabular}{rl}
    \textbf{Versão:} & 1.0 --- Pós-Auditoria de Concorrência \\[0.35em]
    \textbf{Data:} & 23 de Outubro de 2026 \\[0.35em]
    \textbf{Âmbito:} & Backend FastAPI, Frontend Tauri/React \\[0.35em]
    \textbf{Classificação:} & Documento Técnico Interno
  \end{tabular}
  \vspace{1.1cm}
  \begin{minipage}{0.9\textwidth}
    \small\textit{Este documento apresenta o estado atual da plataforma PetroLúmen após a implementação de infraestrutura assíncrona e correção de vulnerabilidades críticas de concorrência. Os dados de mercado e projeções de utilização continuam sujeitos a validação.}
  \end{minipage}
\end{titlepage}

% ------------------------------------------------------------
% SUMÁRIO
% ------------------------------------------------------------
\tableofcontents
\clearpage

% ------------------------------------------------------------
% SUMÁRIO EXECUTIVO
% ------------------------------------------------------------
\section*{Sumário Executivo}
\addcontentsline{toc}{section}{Sumário Executivo}

O setor de engenharia de reservatórios de petróleo e gás debate-se com a fragmentação de ferramentas de software, dependência de folhas de cálculo complexas e a falta de integração entre análise de dados, modelação geológica e simulação de fluxo. O \textbf{PetroLúmen} surge como uma resposta a este problema, oferecendo uma aplicação desktop nativa (Tauri + React) apoiada por um backend de alta performance (Python + FastAPI).

A presente versão do PetroLúmen representa um marco crítico: \textbf{a transição de um protótipo funcional para uma arquitetura preparada para multi-utilizador}. Durante o último ciclo de desenvolvimento, a infraestrutura do backend foi substancialmente reestruturada para eliminar falhas de concorrência. Anteriormente, os motores de simulação (IMPES, AI Predictor, Geostatística) partilhavam instâncias globais, o que causava "state-bleed" (contaminação de dados entre diferentes sessões). 

Com a implementação do \textbf{Session Registry} e de \textbf{Background Jobs} assíncronos, o PetroLúmen isola agora o estado computacional de cada utilizador. Isto permite que múltiplos engenheiros executem simulações e treinos de modelos de Machine Learning simultaneamente, sem risco de corrupção de resultados.

\textbf{O objetivo do produto} é reduzir o tempo de ciclo entre a ingestão de dados de poços (LAS/CSV) e a tomada de decisão técnica, integrando no mesmo ambiente: PVT, Declínio de Produção (DCA), Análise de Pressão (RTA), Simulação Numérica 1D (IMPES) e previsões baseadas em Inteligência Artificial (XGBoost/Prophet).

\subsection*{Decisão Recomendada}

A recomendação para a próxima fase (Fase 4) é focar na \textbf{Internacionalização (i18n)}, refinamento de UX/UI e na preparação da infraestrutura para implementação Cloud (SaaS). A robustez do backend (184 testes passados e cobertura de concorrência) permite que a equipa avance para a validação com utilizadores reais (Beta fechado) sem receio de falhas críticas de arquitetura.

\begin{table}[H]
\centering\small
\begin{tabularx}{\textwidth}{>{\bfseries}p{3.4cm} X}
\toprule
\rowcolor{tableheader}\multicolumn{2}{c}{\textcolor{white}{\textbf{Resumo da Proposta de Valor PetroLúmen}}}\\
\midrule
Produto & Plataforma Desktop (Tauri) e API (FastAPI) para engenharia de reservatórios.\\
Cliente-alvo & Engenheiros de reservatórios, geocientistas e analistas de dados da indústria Oil \& Gas.\\
Diferencial Técnico & Motores de cálculo nativos (Python) integrados com UI moderna (React) e isolamento garantido por \textit{Session Registry}.\\
Estado Atual & Backend refatorado (184 testes a passar), concorrência resolvida, tarefas assíncronas (Background Jobs) implementadas.\\
Próximos Passos & Tradução de interface (i18n), testes E2E com Playwright e validação do modelo de negócio (licenciamento).\\
\bottomrule
\end{tabularx}
\caption{Resumo executivo da arquitetura e estado do PetroLúmen}
\end{table}

\newpage

% ------------------------------------------------------------
% 1. CONTEXTO E JUSTIFICATIVA
% ------------------------------------------------------------
\section{Contexto e Justificativa}

\subsection{O Desafio Operacional na Engenharia de Reservatórios}

A engenharia de reservatórios exige a integração de disciplinas complexas: petrofísica, geologia, mecânica dos fluidos e economia. Historicamente, as equipas dependem de software legado (pesado, muitas vezes baseado em licenças locais) ou de folhas de cálculo fragmentadas. O resultado é um elevado esforço de reconciliação de dados e uma latência elevada entre a análise e a decisão.

\begin{itemize}[leftmargin=*]
    \item \textbf{Análise PVT (Pressure-Volume-Temperature):} Cálculos de Z-Factor e EOS (Peng-Robinson) frequentemente isolados de dados de produção.
    \item \textbf{Simulação Numérica:} Ferramentas de simulação de fluxo (IMPES) exigem setups morosos e carecem de interfaces amigáveis.
    \item \textbf{Previsão (ML):} Modelos de previsão (XGBoost, Prophet) requerem pipelines de dados que os engenheiros não têm tempo para construir.
\end{itemize}

\subsection{A Janela de Oportunidade Estratégica}

A modernização da indústria (Digital Oilfield) exige ferramentas que corram nativamente no \textit{desktop} do engenheiro, mas que tenham a flexibilidade de uma aplicação web moderna. O PetroLúmen resolve esta lacuna ao empacotar o frontend em Rust/Tauri (baixo consumo de memória) e o backend em Python (acesso nativo a SciPy, NumPy, Scikit-Learn).

\begin{table}[H]
\centering
\renewcommand{\arraystretch}{1.3}
\begin{tabularx}{\textwidth}{>{\bfseries}l X}
\toprule
\rowcolor{tableheader}
\multicolumn{2}{c}{\textcolor{white}{\textbf{Forças que Impulsionam a Plataforma PetroLúmen}}} \\
\midrule
Performance Local & Tauri permite binários leves, eliminando a latência de aplicações Electron.\\
Isolamento de Cálculo & O novo \textit{SessionRegistry} garante que simulações massivas não bloqueiam o servidor e não se misturam.\\
Integração de Motores & Unificação de DCA, RTA, Nodal Analysis, IMPES e PVT numa única API documentada (OpenAPI).\\
\bottomrule
\end{tabularx}
\caption{Forças impulsionadoras do projeto PetroLúmen}
\end{table}

% ------------------------------------------------------------
% 2. VISÃO DO PRODUTO E ARQUITETURA
% ------------------------------------------------------------
\section{Arquitetura do Sistema}

O PetroLúmen adota uma arquitetura cliente-servidor desacoplada, otimizada para o ambiente de trabalho do engenheiro.

\subsection{Arquitetura Técnica e Stack}

\begin{table}[H]
\centering
\renewcommand{\arraystretch}{1.3}
\begin{tabularx}{\textwidth}{>{\bfseries}p{3.5cm} X}
\toprule
\rowcolor{tableheader}
\multicolumn{2}{c}{\textcolor{white}{\textbf{Stack Tecnológico do PetroLúmen}}} \\
\midrule
Frontend (Desktop) & Next.js 15, React 19, TypeScript, Zustand (Estado Global), Recharts, Tauri v2.\\
Backend (API) & Python 3.12, FastAPI, Uvicorn, Pydantic, SQLAlchemy 2.0, Alembic.\\
Motores Científicos & NumPy, SciPy, Pandas, Scikit-Learn, XGBoost, Prophet, PyVista/VTK.\\
Base de Dados & SQLite (Local/Desktop) ou PostgreSQL (Produção/Cloud).\\
Segurança e Sessão & JWT Bearer, Bcrypt, SessionRegistry (AsyncIO).\\
\bottomrule
\end{tabularx}
\caption{Stack técnico e infraestrutura do PetroLúmen}
\end{table}

\subsection{O Paradigma do \textit{Session Registry} (Inovação Arquitetural)}

Antes da última refatoração, o PetroLúmen utilizava instâncias globais para os motores de simulação. Isto significava que se o Utilizador A estivesse a correr uma simulação com um grid de 10x10 e o Utilizador B corresse um grid de 20x20, os resultados podiam corromper-se mutuamente.

A introdução do \texttt{SessionRegistry} alterou este paradigma:
\begin{enumerate}
    \item A API FastAPI interceta a chamada e identifica o \texttt{user\_id} (via JWT).
    \item O \texttt{SessionRegistry} (com um TTL e \textit{LRU eviction}) instancia ou recupera um motor de cálculo exclusivo para aquele utilizador (ex: um \texttt{ReservoirSimulator} isolado).
    \item O endpoint injeta o motor via \texttt{Depends(get\_res\_simulator)}, garantindo isolamento absoluto.
\end{enumerate}

\subsection{Processamento Assíncrono (Background Jobs)}

Tarefas de longa duração (como a simulação IMPES ou o treino de modelos XGBoost) deixaram de bloquear o \textit{event loop} do servidor. Foi criada uma tabela de \texttt{Job} na base de dados. O cliente submete a tarefa (HTTP 202 Accepted), recebe um \texttt{job\_id} e pode fazer \textit{polling} ao estado da tarefa, mantendo a interface fluida.

% ------------------------------------------------------------
% 3. WORKFLOW DE UTILIZAÇÃO
% ------------------------------------------------------------
\section{Workflow de Engenharia e Funcionalidades}

\subsection{Fluxo de Trabalho do Utilizador}

\begin{center}
\fbox{\parbox{0.95\textwidth}{
\textbf{Upload de Dados} $\rightarrow$ \textbf{Análise PVT e DCA} $\rightarrow$ \textbf{Setup de Simulação} $\rightarrow$ \textbf{Run Assíncrono} $\rightarrow$ \textbf{Visualização e Exportação (PDF)}
}}
\end{center}

\subsection{Módulos Integrados}

\begin{table}[H]
\centering\small
\begin{tabularx}{\textwidth}{>{\bfseries}p{2.5cm} X X}
\toprule
\rowcolor{tableheader}
\textcolor{white}{\textbf{Módulo}} & \textcolor{white}{\textbf{Descrição Técnica}} & \textcolor{white}{\textbf{Estado}} \\
\midrule
PVT Engine & Cálculos de Z-Factor (DAK, Hall-Yarborough), EOS (Peng-Robinson) e Viscosidades. & Operacional\\
Flow Simulation & Simulador 1D IMPES (Implicit Pressure, Explicit Saturation) com matriz tridiagonal. & Operacional (Assíncrono)\\
DCA (Decline) & Arps (Exponencial, Hiperbólico, Harmónico), Duong, PLE, SEPD. & Operacional (Stateless)\\
RTA (Transient) & Flowing Material Balance (FMB) para gás e análise pseudo-tempo. & Operacional\\
Nodal Analysis & Resolução de interseção IPR/TPR via \texttt{scipy.optimize.fsolve}. & Operacional\\
AI Forecasting & Modelos XGBoost e Prophet para previsão de produção. & Operacional (Isolado)\\
\bottomrule
\end{tabularx}
\caption{Módulos integrados no Backend e o seu estado de maturidade}
\end{table}

% ------------------------------------------------------------
% 4. PRÉ-CONDIÇÕES E RISCOS
% ------------------------------------------------------------
\section{Segurança, Integridade e Riscos}

\subsection{Conformidade e Isolamento de Dados}

A auditoria e subsequente refatoração eliminaram o risco crítico de \textbf{cross-talk} de dados. Com a implementação do \texttt{USE\_SESSION\_REGISTRY=True} (via Pydantic BaseSettings), o sistema assegura que os cálculos de um utilizador estão logicamente separados na RAM do servidor, prevenindo fugas de dados confidenciais (Data Leakage) de diferentes ativos ou operadores.

\subsection{Matriz de Riscos}

\begin{table}[H]
\centering
\renewcommand{\arraystretch}{1.3}
\small
\begin{tabularx}{\textwidth}{>{\bfseries}p{3.5cm} c c X}
\toprule
\rowcolor{tableheader}
\textcolor{white}{\textbf{Risco}} & \textcolor{white}{\textbf{Prob.}} & \textcolor{white}{\textbf{Impacto}} & \textcolor{white}{\textbf{Mitigação}} \\
\midrule
Vazamento de Memória no Registry & Média & Alto & Implementação de \textit{janitor task} com TTL (Time-To-Live) e LRU.\\
Falha de Precisão Numérica & Baixa & Alto & Validação com dados sintéticos (Dry Gas) e cobertura de testes (184 testes).\\
Bloqueio do Event Loop & Média & Médio & Utilização de \texttt{BackgroundTasks} e afastamento do processamento pesado do endpoint principal.\\
Incompatibilidade de Ficheiros & Média & Médio & Normalização com Pandas e suporte estrito a formatos LAS/CSV no uploader.\\
\bottomrule
\end{tabularx}
\caption{Matriz de riscos técnicos do PetroLúmen (Pós-Fase 2)}
\end{table}

% ------------------------------------------------------------
% 5. ROADMAP E PLANO DE IMPLEMENTAÇÃO
% ------------------------------------------------------------
\section{Roadmap e Próximas Fases}

O projeto encontra-se numa fase de maturidade técnica elevada. A transição para a Fase 4 foca-se na expansão da experiência do utilizador e na preparação comercial.

\begin{table}[H]
\centering\small
\begin{tabularx}{\textwidth}{>{\bfseries}p{2cm} p{4cm} X}
\toprule
\rowcolor{tableheader}\textcolor{white}{Fase} & \textcolor{white}{Atividade Principal} & \textcolor{white}{Entregável}\\
\midrule
Concluída & Concorrência e Background Jobs & Session Registry, Endpoints Assíncronos, 184 testes a passar.\\
Fase 4.1 & i18n e L10n & Migração do frontend para \texttt{next-intl} e padronização PT/EN.\\
Fase 4.2 & E2E Testing & Criação de testes Playwright para o fluxo completo (Upload $\rightarrow$ Simulação $\rightarrow$ PDF).\\
Fase 4.3 & Empacotamento e Release & Geração do instalador final Tauri (.msi, .dmg, .AppImage) com auto-updater.\\
Fase 5.0 & Validação com Utilizadores & Beta fechado com engenheiros reais para testar precisão e UX.\\
\bottomrule
\end{tabularx}
\caption{Plano de transição para a versão comercial}
\end{table}

% ------------------------------------------------------------
% 6. CONCLUSÃO
% ------------------------------------------------------------
\section{Conclusão}

O PetroLúmen deixou de ser um conjunto de scripts de engenharia para se tornar uma plataforma de software robusta, segura e escalável. A erradicação de falhas de concorrência (através do \textit{Session Registry}) e a implementação de tarefas assíncronas representam um marco de engenharia de software que valida a viabilidade técnica do projeto.

A plataforma está agora pronta para ser testada em ambientes mais próximos da produção (Beta), onde o isolamento de dados, a estabilidade numérica e a resposta rápida da interface serão os principais indicadores de sucesso.

\vspace{0.5cm}
\noindent\fbox{\parbox{\dimexpr\textwidth-2\fboxsep-2\fboxrule}{%
\textbf{Decisão Proposta:} Avançar para a Fase 4 (i18n, Testes E2E e Empacotamento Final) e iniciar o recrutamento de um grupo restrito de engenheiros de reservatórios para o programa Beta.}}

\vspace{1cm}

% ------------------------------------------------------------
% GLOSSÁRIO E ANEXOS
% ------------------------------------------------------------
\newpage
\section*{Glossário Técnico}
\addcontentsline{toc}{section}{Glossário Técnico}

\begin{description}[leftmargin=1em, style=nextline]
    \item[API] \textit{Application Programming Interface} --- Interface que permite a comunicação entre o Frontend e o Backend.
    \item[DCA] \textit{Decline Curve Analysis} --- Análise de Curva de Declínio de produção.
    \item[EOS] \textit{Equation of State} --- Equação de Estado (ex: Peng-Robinson) usada para modelar o comportamento de fluidos.
    \item[IMPES] \textit{Implicit Pressure, Explicit Saturation} --- Método numérico para simulação de reservatórios.
    \item[JWT] \textit{JSON Web Token} --- Padrão seguro para transmissão de autenticação entre cliente e servidor.
    \item[LRU] \textit{Least Recently Used} --- Política de cache usada no \textit{Session Registry} para eliminar sessões antigas.
    \item[PVT] \textit{Pressure-Volume-Temperature} --- Propriedades termodinâmicas dos fluidos do reservatório.
    \item[RTA] \textit{Rate Transient Analysis} --- Análise de dados de pressão e vazão em regime transiente.
    \item[Tauri] --- \textit{Framework} em Rust usado para empacotar a aplicação web num executável nativo de desktop.
    \item[TTL] \textit{Time-To-Live} --- Tempo de vida útil de uma sessão antes de ser destruída na RAM.
\end{description}

\end{document}